(1) 线性连续方程为 \(\eta_t+h u_x=0\)，动量方程 \(u_t+g\eta_x=0\)，合并给 \(\eta_{tt}=gh\,\eta_{xx}\)，故相速 \(c=\sqrt{gh}\)。对于向右行水波，速度与高程满足 \(u=(c/h)\eta=\sqrt{g/h}\eta\)；向左行波则 \(u=-(c/h)\eta\)。这种符号区别是反射界面质量通量条件的关键。

(2) 写入射、反射、透射水位振幅分别为 \(A_i,A_r,A_t\)。界面水位连续 \(A_i+A_r=A_t\)，体积通量连续 \(h_1u_1=h_2u_2\)，给 \(c_1(A_i-A_r)=c_2A_t\)。求得水位振幅反射系数 \(r=(c_1-c_2)/(c_1+c_2)\)，透射系数 \(t=2c_1/(c_1+c_2)\)。若 \(c_1=c_2\)，反射消失；注意不同水深的水位振幅不同于速度振幅。

(3) 单位宽水波平均能量密度为 \(\rho g|A|^2/2\)，能流大小为此密度乘相速 \(c\)，所以反射率 \(R=|r|^2\)，透射率 \(\mathcal T=(c_2/c_1)|t|^2=4c_1c_2/(c_1+c_2)^2\)。直接相加得 \((c_1-c_2)^2/(c_1+c_2)^2+4c_1c_2/(c_1+c_2)^2=1\)。若错误地用 \(t^2\) 作透射率，将违反能量守恒。

(4) 当从深水到浅水 \(h_2<h_1\) 时 \(c_2<c_1\)，水位透射振幅 \(t>1\)，意味着浅水一侧波峰可高于入射波峰，但能量透射率仍不超过一。反过来浅到深时 \(r<0\) 表示反射水位反相。在极端浅水区域线性浅水假设最终会因振幅相对水深过大而失效，实际海啸会涉及非线性、底部摩擦和碎波，而不能将振幅增长无界外推。
